\section{Discussion}
\label{sec:discussion}

\subsection{Embedding-Model Performance}

The superiority of 1024-dimension models (E5-large, BGE-M3) over 768-dimension
models confirms the importance of representational capacity for Arabic legal
vocabulary. The extra 256 dimensions encode fine semantic nuances crucial for
distinguishing close articles such as simple theft (Art.~147) from aggravated
theft (Art.~148) in the Moroccan Penal Code. The gap between E5-large and
BGE-M3 is not statistically significant ($p = 0.08$), making them
interchangeable for production. The small performance drop for cross-lingual
queries (French MRR 0.85 vs.\ formal Arabic MRR 0.94) shows that modern
multilingual embedding models have reached sufficient cross-lingual
generalization for specialized applications.

\subsection{BM25 Contribution}

With the specialized Arabic tokenizer, BM25 now achieves non-zero scores
(Recall@5 = 41.67\%, Hit@5 = 46.88\%), unlike the previous PostgreSQL FTS
implementation that failed entirely. BM25 is particularly useful for
article-number lookups and specific legal-term matching. In RRF, BM25
contributes complementary results, improving Hit@5 from 65.62\% (dense) to
68.75\% (RRF) and bringing RRF's MRR@5 (0.533) close to Dense's (0.565).

\subsection{Faithfulness and Answer Quality}

The measured faithfulness (0.47 for dense) should not be interpreted solely as
a hallucination rate. Qualitative analysis reveals that the LLM produces
legitimate contextual enrichment -- multi-article synthesis, structured
citations, and reformulation of legislative language into accessible
explanations -- that token-overlap metrics fail to capture. Recent surveys of hallucination
in natural language generation~\cite{ji2023survey} similarly emphasize the
need for multi-faceted faithfulness evaluation, and an
LLM-as-judge evaluation~\cite{gao2023faithful} would provide a more accurate
assessment. Answer relevance (0.69) is satisfactory, indicating that the LLM
correctly addresses the question in most cases.

\subsection{Limitations}

\begin{enumerate}
  \item \textbf{Darija--MSA gap}: a 29\% relative performance drop confirms
        Arabic diglossia as an open challenge.
  \item \textbf{Tax-domain failure}: all metrics are zero due to a domain
        taxonomy mismatch (articles stored as \texttt{civil}, queried as
        \texttt{tax}).
  \item \textbf{Latency}: 13--14 seconds total, bottlenecked by the remote
        Cloudflare-tunnel embedding service.
  \item \textbf{Corpus coverage}: limited to 7 legal domains; full coverage
        of Moroccan law requires additional computational resources.
\end{enumerate}

\subsection{Future Work}

Several directions are planned: (1) embedding fine-tuning on the Moroccan
legal corpus to reduce the Darija--MSA gap; (2) local embedding deployment to
reduce latency from $\sim$5.5\,s to $\sim$100\,ms; (3) LLM-as-judge evaluation
for more accurate faithfulness measurement; (4) extension of the corpus to
cover the complete Moroccan Official Bulletin; (5) multi-turn conversational
support with context memory; and (6) an upgrade to Gemma 4 9B (or equivalent)
for improved generation quality.
